\documentclass[11pt]{article}

\usepackage[a4paper,margin=1in]{geometry}
\usepackage{amsmath,amssymb,bm}
\usepackage{booktabs}
\usepackage{hyperref}

\hypersetup{
    colorlinks=true,
    linkcolor=blue,
    citecolor=blue,
    urlcolor=blue
}

\title{Toward an Averaged GEqOE Formulation for Autonomous Zonal Dynamics\\
\large \texorpdfstring{$J_2$}{J2} Closure, General Zonal Structure, and a Path to Fourier-in-\texorpdfstring{$K$}{K} Thrust Averaging}
\author{ASTRODYN-CORE Internal Research Note}
\date{\today}

\begin{document}
\maketitle

\begin{abstract}
This note specializes the mean / averaged GEqOE question to the autonomous
conservative case generated by the Earth zonal harmonics. The objective is to
identify a tractable research path toward a fast analytical or semi-analytical
model that could later support maneuver-detection metrics.

Three results are recorded. First, for a static zonal potential, the current
GEqOE equations already admit an exact first-order orbit-averaging operator in
the propagated generalized eccentric longitude $K$. Second, for $J_2$ alone,
the first-order secular closure is especially simple: $\nu$, the generalized
eccentricity magnitude $g=\sqrt{p_1^2+p_2^2}$, and the inclination magnitude
$Q=\sqrt{q_1^2+q_2^2}$ are constants, while the pairs $(p_1,p_2)$ and
$(q_1,q_2)$ undergo uniform rotations with explicit secular rates. Third, the
same $K$-based framework provides a clean insertion point for a future
Fourier-in-$K$ thrust averaging theory, provided the first step is taken around
the autonomous $J_2$ or zonal baseline.
\end{abstract}

\section{Scope}

The present note focuses on the GEqOE state propagated in the current branch:
\begin{equation}
\bm{y} = (\nu,p_1,p_2,K,q_1,q_2).
\end{equation}
The working assumptions are:
\begin{enumerate}
    \item the conservative perturbation is autonomous and axisymmetric about the
    inertial $z$ axis;
    \item the disturbing potential is either the pure $J_2$ term or a static
    zonal expansion;
    \item averaging is performed with respect to the propagated fast angle $K$;
    \item only first-order secular dynamics are sought at this stage.
\end{enumerate}

This restriction is deliberate. It keeps the averaged problem autonomous and
avoids the additional complications introduced by explicit time dependence
(third-body ephemerides, rotating-frame geopotential tides, etc.).

\section{GEqOE Preliminaries}

Define the auxiliary quantities
\begin{align}
g^2 &= p_1^2 + p_2^2, \qquad \beta = \sqrt{1-g^2}, \qquad
\alpha = \frac{1}{1+\beta}, \\
a &= \left(\frac{\mu}{\nu^2}\right)^{1/3}, \\
Q^2 &= q_1^2 + q_2^2, \qquad
\gamma = 1 + Q^2, \qquad
\delta = 1 - Q^2.
\end{align}
From the equinoctial orientation variables one has
\begin{equation}
\cos i = \frac{\delta}{\gamma},
\qquad
\tan\frac{i}{2} = Q.
\end{equation}

The propagated generalized eccentric longitude $K$ gives
\begin{align}
r &= a\left(1 - p_1 \sin K - p_2 \cos K\right), \label{eq:rK}\\
X &= a\left(\alpha p_1 p_2 \sin K + (1-\alpha p_1^2)\cos K - p_2\right), \\
Y &= a\left(\alpha p_1 p_2 \cos K + (1-\alpha p_2^2)\sin K - p_1\right),
\end{align}
and the normalized inertial $z$ coordinate is
\begin{equation}
\hat z = \frac{2(Yq_2 - Xq_1)}{\gamma r}. \label{eq:zhat}
\end{equation}

The generalized angular momentum proxy used throughout the GEqOE paper and the
current implementation is
\begin{equation}
c = \left(\frac{\mu^2}{\nu}\right)^{1/3}\beta. \label{eq:cdef}
\end{equation}
Under purely conservative autonomous dynamics, the generalized energy is a first
integral and therefore
\begin{equation}
\dot{\nu} = 0.
\end{equation}

\section{Autonomous Zonal GEqOE Dynamics}

For a zonal potential of the form
\begin{equation}
U(r,\hat z) = \sum_{n\ge 2} U_n(r,\hat z)
= \sum_{n\ge 2} \frac{C_n}{r^{n+1}} P_n(\hat z),
\qquad
C_n = \mu J_n R_e^n,
\end{equation}
the current branch uses the exact conservative GEqOE system
\begin{align}
\dot{\nu} &= 0, \\
\dot{p}_1 &= p_2(d-w_h)
+ \frac{1}{c}\left(\frac{X}{a}+2p_2\right)\mathcal{E}_Z, \label{eq:p1_zonal}\\
\dot{p}_2 &= p_1(w_h-d)
- \frac{1}{c}\left(\frac{Y}{a}+2p_1\right)\mathcal{E}_Z, \label{eq:p2_zonal}\\
\dot{K} &= \frac{w}{r} + d - w_h
+ \frac{1}{c}\left(1+\alpha\left(1-\frac{r}{a}\right)\right)\mathcal{E}_Z, \\
\dot{q}_1 &= \frac{\gamma}{2} w_Y, \label{eq:q1_zonal}\\
\dot{q}_2 &= \frac{\gamma}{2} w_X, \label{eq:q2_zonal}
\end{align}
with
\begin{align}
w &= \sqrt{\frac{\mu}{a}}, \\
h &= \sqrt{c^2 - 2r^2 U}, \\
d &= \frac{h-c}{r^2}, \\
\mathcal{E}_Z &= \sum_{n\ge 2} (1-n) U_n, \\
F_h &= - \frac{\delta}{\gamma r}\,\frac{\partial U}{\partial \hat z}, \\
w_X &= \frac{X}{h}F_h, \qquad
w_Y = \frac{Y}{h}F_h, \qquad
w_h = w_X q_1 - w_Y q_2.
\end{align}

This formulation is already structurally suitable for averaging because:
\begin{enumerate}
    \item the system is autonomous;
    \item the fast angle is explicit;
    \item $\dot{\nu}=0$ exactly;
    \item all forcing enters through $r(K)$, $X(K)$, $Y(K)$, and $\hat z(K)$.
\end{enumerate}

\section{First-Order Orbit Averaging in \texorpdfstring{$K$}{K}}

\subsection{Averaging operator}

To first order in the zonal perturbation, the unperturbed fast flow is
\begin{equation}
\dot{K}_0 = \frac{w}{r},
\qquad
\frac{dt}{dK} = \frac{r}{w}.
\end{equation}
Hence the orbital average of any $2\pi$-periodic quantity $f(K)$ is
\begin{equation}
\langle f \rangle_K
= \frac{1}{T}\int_0^T f\,dt
= \frac{1}{2\pi a}\int_0^{2\pi} r(K)\,f(K)\,dK.
\label{eq:avg_operator}
\end{equation}

Let the slow state be
\begin{equation}
\bm{z} = (\nu,p_1,p_2,q_1,q_2).
\end{equation}
Then the exact first-order averaged slow flow of the autonomous zonal problem is
\begin{align}
\dot{\bar\nu} &= 0, \\
\dot{\bar p}_1 &=
\frac{1}{2\pi \bar a}\int_0^{2\pi} \bar r
\left[
\bar p_2(\bar d-\bar w_h)
+ \frac{1}{\bar c}\left(\frac{\bar X}{\bar a}+2\bar p_2\right)\bar{\mathcal E}_Z
\right]\,dK, \label{eq:avg_p1_general}\\
\dot{\bar p}_2 &=
\frac{1}{2\pi \bar a}\int_0^{2\pi} \bar r
\left[
\bar p_1(\bar w_h-\bar d)
- \frac{1}{\bar c}\left(\frac{\bar Y}{\bar a}+2\bar p_1\right)\bar{\mathcal E}_Z
\right]\,dK, \label{eq:avg_p2_general}\\
\dot{\bar q}_1 &=
\frac{1}{2\pi \bar a}\int_0^{2\pi}
\bar r \,\frac{\bar\gamma}{2}\,\bar w_Y\,dK, \label{eq:avg_q1_general}\\
\dot{\bar q}_2 &=
\frac{1}{2\pi \bar a}\int_0^{2\pi}
\bar r \,\frac{\bar\gamma}{2}\,\bar w_X\,dK. \label{eq:avg_q2_general}
\end{align}
All barred intermediate quantities are evaluated from the frozen slow state
$\bar{\bm z}$ and the fast phase $K$ via the unperturbed relations
\eqref{eq:rK}-\eqref{eq:zhat} and \eqref{eq:cdef}.

\subsection{Interpretation}

Equations \eqref{eq:avg_p1_general}-\eqref{eq:avg_q2_general} are already a
valid first-order mean-element model for any autonomous zonal expansion.
They are not yet a closed analytical solution, but they are an exact
one-revolution quadrature formula in the variables used by the present branch.

This is the right point of departure for research, because it separates the
problem into two layers:
\begin{enumerate}
    \item a \emph{formal averaging layer}, which is already solved by
    \eqref{eq:avg_operator}-\eqref{eq:avg_q2_general};
    \item a \emph{closure layer}, where one asks whether the integrals can be
    reduced to explicit analytical expressions.
\end{enumerate}

\section{Closed First-Order \texorpdfstring{$J_2$}{J2} Secular Model}

\subsection{Why \texorpdfstring{$J_2$}{J2} is special}

For the pure $J_2$ potential,
\begin{equation}
U_{J_2} = -\frac{A}{r^3}\left(1-3\hat z^2\right),
\qquad
A = \frac{\mu J_2 R_e^2}{2},
\end{equation}
the averaged disturbing function is exceptionally simple. Using the standard
Keplerian averages
\begin{equation}
\left\langle \frac{1}{r^3}\right\rangle
= \frac{1}{a^3\beta^3},
\qquad
\left\langle \frac{\sin^2 u}{r^3}\right\rangle
= \frac{1}{2a^3\beta^3},
\end{equation}
with $u=\omega+f$, one obtains
\begin{equation}
\bar U_{J_2}
= -\frac{A}{2a^3\beta^3}\left(3\cos^2 i - 1\right)
= -\frac{A}{2a^3\beta^3}
\left(3\frac{\delta^2}{\gamma^2} - 1\right).
\label{eq:Ubar_J2}
\end{equation}

Equation \eqref{eq:Ubar_J2} depends only on $a$, $\beta$, and $i$. It has no
dependence on the node or pericenter angles. Consequently, at first order in
$J_2$,
\begin{equation}
\dot{\bar\nu} = 0,
\qquad
\dot{\bar g} = 0,
\qquad
\dot{\bar Q} = 0.
\end{equation}
Thus the magnitudes of $(p_1,p_2)$ and $(q_1,q_2)$ are constant and the secular
dynamics reduce to pure rotations.

\subsection{Secular rates}

Let
\begin{equation}
c_i = \cos i = \frac{\delta}{\gamma},
\qquad
p = a\beta^2.
\end{equation}
At first order in $J_2$, the secular nodal rate is
\begin{equation}
\omega_{\Omega,J_2}
= -\frac{3}{2}\,\nu J_2
\left(\frac{R_e}{p}\right)^2 c_i
= -\frac{3A}{c a^3 \beta^3}\,\frac{\delta}{\gamma},
\label{eq:OmegaJ2}
\end{equation}
and the secular rate of the generalized longitude of pericenter is
\begin{equation}
\omega_{\Psi,J_2}
= \frac{3}{4}\,\nu J_2
\left(\frac{R_e}{p}\right)^2
\left(5c_i^2 - 2c_i - 1\right)
= \frac{3A}{2c a^3 \beta^3}
\left(
5\frac{\delta^2}{\gamma^2}
- 2\frac{\delta}{\gamma}
- 1
\right).
\label{eq:PsiJ2}
\end{equation}

Because
\begin{equation}
p_1 = g\sin\Psi,
\qquad
p_2 = g\cos\Psi,
\qquad
q_1 = Q\sin\Omega,
\qquad
q_2 = Q\cos\Omega,
\end{equation}
the first-order mean GEqOE system for $J_2$ becomes
\begin{align}
\dot{\bar\nu} &= 0, \label{eq:j2_mean_nu}\\
\dot{\bar p}_1 &= \omega_{\Psi,J_2}\,\bar p_2, \label{eq:j2_mean_p1}\\
\dot{\bar p}_2 &= -\omega_{\Psi,J_2}\,\bar p_1, \label{eq:j2_mean_p2}\\
\dot{\bar q}_1 &= \omega_{\Omega,J_2}\,\bar q_2, \label{eq:j2_mean_q1}\\
\dot{\bar q}_2 &= -\omega_{\Omega,J_2}\,\bar q_1. \label{eq:j2_mean_q2}
\end{align}

Equations \eqref{eq:j2_mean_nu}-\eqref{eq:j2_mean_q2} are autonomous and closed.
They provide exactly the kind of reduced slow-flow model that is suitable for
fast metric computation.

\subsection{Mean longitude / phase}

If a mean fast phase is also required, the natural choice is the mean
generalized longitude
\begin{equation}
\bar L = \bar M + \bar\Psi.
\end{equation}
Using the standard first-order $J_2$ secular rate of the mean anomaly,
\begin{equation}
\dot{\bar M}
= \nu
+ \frac{3}{4}\,\nu J_2
\left(\frac{R_e}{a\beta^2}\right)^2
\beta\left(3c_i^2-1\right),
\end{equation}
one obtains
\begin{equation}
\dot{\bar L}
= \nu
+ \frac{3}{4}\,\nu J_2
\left(\frac{R_e}{a\beta^2}\right)^2
\left[
\beta(3c_i^2-1)
+ 5c_i^2 - 2c_i - 1
\right].
\label{eq:Lbar_J2}
\end{equation}
For many maneuver-detection applications, however, the autonomous slow
subsystem \eqref{eq:j2_mean_nu}-\eqref{eq:j2_mean_q2} is the more useful core
object, while the phase can be reconstructed only when needed.

\section{What Changes for General Zonal Expansions}

\subsection{Autonomy is preserved}

For a general static zonal potential, the averaged GEqOE system remains
autonomous. This is the main reason to attack the zonal problem before adding
non-autonomous perturbations. The exact first-order mean system is already given
by \eqref{eq:avg_p1_general}-\eqref{eq:avg_q2_general}.

\subsection{\texorpdfstring{$J_2$}{J2} is not representative of all zonals}

The $J_2$ closure is unusually simple because the averaged disturbing function
depends only on $a$, $e$, and $i$. Higher zonals are not guaranteed to share
that property:
\begin{enumerate}
    \item odd zonals such as $J_3$ generally retain periapsis dependence in the
    singly averaged problem;
    \item higher even zonals can also introduce more complicated slow-angle
    structure than the pure $J_2$ case;
    \item the slow subsystem therefore need not remain a uniform rotation in
    $(p_1,p_2)$.
\end{enumerate}

What does survive is autonomy and axisymmetry. Hence the averaged vector field
cannot depend on absolute time and cannot depend on the node angle by itself.
In equinoctial variables this implies that the mean zonal dynamics can only
depend on:
\begin{enumerate}
    \item the invariants $a$, $g$, and $Q$;
    \item the relative pericenter-orientation information carried by
    $(p_1,p_2)$ with respect to the symmetry axis encoded by $(q_1,q_2)$.
\end{enumerate}

\subsection{Recommended next derivation target}

The next mathematically useful target is not ``all zonals at once'' but:
\begin{enumerate}
    \item derive the closed first-order $J_2$ system rigorously in GEqOE
    notation, validating \eqref{eq:j2_mean_nu}-\eqref{eq:j2_mean_q2};
    \item then evaluate the quadratures
    \eqref{eq:avg_p1_general}-\eqref{eq:avg_q2_general} term by term for
    $J_2+J_3+J_4$;
    \item then decide whether a compact analytical closure exists or whether a
    semi-analytical quadrature is the more pragmatic model.
\end{enumerate}

\section{Path to a Fourier-in-\texorpdfstring{$K$}{K} Thrust Extension}

Consider adding a small non-conservative acceleration expressed in the orbital
RTN frame,
\begin{equation}
\bm{P}(K)
= P_r(K)\,\bm{e}_r + P_t(K)\,\bm{e}_t + P_h(K)\,\bm{e}_h,
\end{equation}
with Fourier expansions in the propagated GEqOE fast angle $K$,
\begin{align}
P_r(K) &= a_{r0} + \sum_{m=1}^M \left(a_{rm}\cos mK + b_{rm}\sin mK\right), \\
P_t(K) &= a_{t0} + \sum_{m=1}^M \left(a_{tm}\cos mK + b_{tm}\sin mK\right), \\
P_h(K) &= a_{h0} + \sum_{m=1}^M \left(a_{hm}\cos mK + b_{hm}\sin mK\right).
\end{align}

Let $\bm{c}$ denote the stacked Fourier coefficients. If one linearizes the
forced averaged problem about the autonomous $J_2$ slow flow and neglects the
mixed $J_2$-thrust coupling at first pass, the mean dynamics take the form
\begin{equation}
\dot{\bar{\bm z}}
= \bm{f}^{\,sec}_{J_2}(\bar{\bm z})
+ \bm{B}(\bar{\bm z})\,\bm{c}
+ O(J_2\|\bm{c}\|,\|\bm{c}\|^2),
\label{eq:forced_avg_form}
\end{equation}
where $\bm{f}^{\,sec}_{J_2}$ is the autonomous slow flow
\eqref{eq:j2_mean_nu}-\eqref{eq:j2_mean_q2}.

The main structural reason this is promising is that in GEqOE:
\begin{enumerate}
    \item $r/a$ is affine in $\sin K$ and $\cos K$;
    \item $X/a$ and $Y/a$ are affine in $\sin K$ and $\cos K$;
    \item the averaging weight is also $r/a$.
\end{enumerate}
Therefore the thrust-induced averaged kernels are low-order trigonometric
polynomials in $K$. Before retaining $J_2$-thrust coupling, the resulting
matrix $\bm{B}(\bar{\bm z})$ should depend only on a finite low-order subset of
the Fourier coefficients. This is the GEqOE analogue of the classical thrust
Fourier coefficient reduction, and it is the correct next target once the
autonomous $J_2$ baseline is fully controlled.

\section{Conclusions}

For the autonomous zonal problem, the averaged GEqOE program is viable.

\begin{enumerate}
    \item The current branch already contains the right dynamical variables and
    the right fast angle for averaging: the generalized eccentric longitude $K$.
    \item For a general static zonal potential, the first-order mean GEqOE
    equations exist immediately as the autonomous quadratures
    \eqref{eq:avg_p1_general}-\eqref{eq:avg_q2_general}.
    \item For pure $J_2$, a compact autonomous secular closure is available:
    $\nu$, $g$, and $Q$ are constant, while $(p_1,p_2)$ and $(q_1,q_2)$ rotate
    with the explicit rates \eqref{eq:OmegaJ2}-\eqref{eq:PsiJ2}.
    \item This $J_2$ closure is a credible baseline for fast maneuver metrics.
    \item The natural next extension is a Fourier-in-$K$ thrust perturbation on
    top of the autonomous $J_2$ mean flow, not a jump directly to fully general
    non-autonomous perturbations.
\end{enumerate}

In short: the autonomous $J_2$ problem is the correct place to start, the
general zonal problem is still structurally tractable, and the thrust-Fourier
extension should be built on top of this averaged baseline rather than beside
it.

\end{document}
